% Generated draft: requires mathematical and source review.
\documentclass[11pt]{article}
\usepackage[margin=25mm]{geometry}
\usepackage{fontspec}
\setmainfont{Latin Modern Roman}
\usepackage{amsmath,amssymb,mathtools,unicode-math}
\setmathfont{Latin Modern Math}
\usepackage{enumitem,needspace,xurl,hyperref}
\setlength{\parindent}{0pt}
\setlength{\parskip}{5pt}
\setlength{\emergencystretch}{3em}
\newcommand{\N}{\mathbb{N}}
\newcommand{\Z}{\mathbb{Z}}
\newcommand{\Q}{\mathbb{Q}}
\newcommand{\R}{\mathbb{R}}
\newcommand{\C}{\mathbb{C}}
\newcommand{\F}{\mathbb{F}}
\newcommand{\A}{\mathbb{A}}
\newcommand{\PP}{\mathbb{P}}
\newcommand{\E}{\mathbb{E}}
\newcommand{\eps}{\varepsilon}
\newcommand{\dd}{\,\mathrm{d}}
\begin{document}
\section*{UnsolvedMath ID: 30001630}
\textbf{Source problem code:} catalogue entry
\subsection{Definitions and mathematical statement}
A compact surface is a surface whose underlying topological space is compact. A puncture is a removed point of a surface. These terms alone do not specify a geometric structure, an operator, or an index theorem.

\textbf{Source question.} Can the index theorem for punctured surfaces be derived from the index theorem for compact surfaces through degeneration?

The intended task is to obtain the punctured-surface index theorem by a justified degeneration argument starting from the compact-surface index theorem. The supplied record does not identify either theorem or specify what constitutes such a derivation. In particular, it does not specify whether the theorem concerns a single operator or a family, the class of surfaces, the number of punctures, metrics or other geometric data, operator domains, conditions at punctures or boundaries, or the degeneration and its limiting terms. Consequently, a precise quantified theorem cannot be extracted from this record alone. These choices remain unresolved; no particular index formula or degeneration model is asserted here.
\subsection{Short English statement}
The question asks whether degeneration of compact surfaces can yield an index theorem for punctured surfaces. The intended index theorems and the hypotheses governing the degeneration are unspecified, so this remains a statement-only draft requiring mathematical review.
\subsection{Sources}
\begin{itemize}
\item Supplied local UnsolvedMath registry record, permanent ID \(30001630\), source problem code \texttt{catalogue entry}, title ``Degeneration Proofs of Index Theorems for Punctured Surfaces,'' in the collection ``Oberwolfach Reports Open Problems.'' The record gives proposed year \(2010\) and identifies no individual proposer.
\item Supplied source citation: \emph{Teichmüller Theory (2011)}, Oberwolfach Reports, DOI \texttt{10.4171/owr/2010/53}, pp.\ 3085--3157. \url{https://doi.org/10.4171/owr/2010/53}. The background attributes the question to workshop open-problem discussions documented in Oberwolfach Reports.
\item Supplied literature source, labeled ``original OWR report'': \url{https://doi.org/10.4171/owr/2010/53}.
\item Supplied literature source, labeled ``analytic torsion for surfaces with cusps'': \url{https://arxiv.org/abs/1812.11739}.
\item Literature assessment dated exactly \texttt{2026-08-21}: the supplied curation labels the problem ``partially solved'' and states that index and curvature formulas for punctured families have been substantially generalized, including analytic-torsion and cusp-curvature results, while the specific degeneration derivation from the compact index theorem is not established in full generality. This is supplied dated curation, not an independently verified literature assessment; its time-sensitive status and cited claims require independent verification.
\item The supplied record labels the statement ``exact'' and reports that its mathematical content was checked against the official Oberwolfach report, with notation, context, and question form normalized without changing the source-level claim. This is a supplied verification claim, not an independent source check performed for this draft.
\end{itemize}
\end{document}
